%% ----------------------------------------------------------------------------
%%  §3.9  Appendice: infiniti di ordine crescente
%% ----------------------------------------------------------------------------
\section{Appendice: infiniti di ordine crescente}
\label{sec:infiniti}

Molte successioni importanti divergono a $+\infty$, ma con «velocità» molto diverse. In
questo paragrafo impariamo a confrontarle; lo strumento principale è un criterio che
sfrutta il teorema sulle successioni monotòne, con lo stesso metodo utilizzato per il
metodo di Erone.

\begin{teorema}[Criterio del rapporto per le successioni]\label{teo:criterio-rapporto}\index{criterio del rapporto!per le successioni}
Sia $a_{n}$ una successione a termini positivi. Definiamo $b_{n}=a_{n+1}/a_{n}$. Se la
successione $b_{n}$ converge a un limite $b<1$, allora la successione $a_{n}$ tende a
zero.
\end{teorema}

\begin{schema}
\Ipotesi{$a_{n}>0$ per ogni $n$; \ $\dfrac{a_{n+1}}{a_{n}}\to b$, con $b<1$.}
\Tesi{$a_{n}\to0$.}
\end{schema}

\begin{strategia}
Se il rapporto tra un termine e il precedente tende a un numero minore di $1$, da un
certo punto in poi ogni termine è più piccolo del precedente: la successione è
definitivamente decrescente e, essendo positiva, limitata inferiormente da $0$; quindi
converge (\cref{teo:monotone}) a un limite $a\ge0$. Per identificare $a$ si passa al
limite nella relazione $a_{n+1}=b_{n}a_{n}$, come nel metodo di Erone: se fosse $a>0$, si
otterrebbe $b=a/a=1$, contro l'ipotesi.
\end{strategia}

\begin{dimostrazione}
Per il teorema della permanenza del segno (\cref{teo:permanenza-segno}), applicato alla
successione $1-b_{n}$, che converge a $1-b>0$, esiste un indice $\nu$ per cui $b_{n}<1$
per ogni $n>\nu$. Quindi $a_{n+1}/a_{n}<1$, cioè $a_{n+1}<a_{n}$, per ogni $n>\nu$. Il
teorema sulle successioni monotòne (nella forma della \cref{oss:traslata}) assicura
l'esistenza del limite $a$, che è un numero reale non negativo, dato che la successione
è definitivamente decrescente e a termini positivi (\cref{cor:limite-non-negativo}).
Supponendo per assurdo $a\neq0$ e passando al limite per $n\to+\infty$ nella relazione
$b_{n}=a_{n+1}/a_{n}$ (si ricordi che $a_{n+1}\to a$, \cref{oss:traslata}), si ottiene
$b=a/a=1$, in contrasto con l'ipotesi $b<1$. Pertanto risulta $a=0$.
\end{dimostrazione}

\begin{osservazione}
Se invece $b>1$, lo stesso ragionamento mostra che $a_{n}$ è definitivamente crescente,
diciamo per $n\ge n_{0}$, e quindi ha limite. Tale limite non può essere $0$, perché è
maggiore o uguale di $a_{n_{0}}>0$, e non può essere un numero positivo $a$, perché si
otterrebbe di nuovo $b=a/a=1$: dunque $a_{n}\to+\infty$. Allo stesso risultato si arriva
applicando il criterio alla successione $1/a_{n}$, il cui rapporto
$\frac{1/a_{n+1}}{1/a_{n}}=\frac1{b_{n}}$ tende a $\frac1b<1$: quindi $1/a_{n}\to0$ e,
poiché $a_{n}>0$, $a_{n}\to+\infty$ (si veda il riquadro «$1/0$ non è $+\infty$» nel
\cref{sec:forme-indeterminate}). È la versione «per divergenza» del criterio. Se infine $b=1$, il criterio non dà alcuna
informazione: $a_{n}=n$, $a_{n}=1$ e $a_{n}=1/n$ hanno tutte $b_{n}\to1$, ma
comportamenti diversi.
\end{osservazione}

\subsection{La gerarchia degli infiniti}

Applichiamo il criterio del rapporto al confronto delle successioni
\begin{equation}\label{eq:gerarchia}
  \log n; \qquad n^{b}; \qquad a^{n}; \qquad n!; \qquad n^{n},
\end{equation}
dove abbiamo scelto $b>0$ e $a>1$. Il simbolo $n!$ ($n$ \emph{fattoriale}) indica il
prodotto dei primi $n$ numeri naturali (\cref{sec:combinatoria}). Per $n\to+\infty$ le
cinque successioni \eqref{eq:gerarchia} tendono tutte a $+\infty$ (per $\log n$ e
$n^{b}$ lo verificheremo nel \cref{lem:potenze-log}, per $a^{n}$ è il
\cref{teo:lim-geometrica}, per $n!\ge n$ e $n^{n}\ge n$ è ovvio). Possiamo però dire che
sono \emph{infiniti di ordine crescente}\index{infiniti di ordine crescente}, nel senso
che i limiti dei rapporti valgono
\begin{equation}\label{eq:gerarchia-limiti}
  \lim_{n\to+\infty}\frac{\log n}{n^{b}}=\lim_{n\to+\infty}\frac{n^{b}}{a^{n}}
  =\lim_{n\to+\infty}\frac{a^{n}}{n!}=\lim_{n\to+\infty}\frac{n!}{n^{n}}=0 .
\end{equation}

\begin{lemma}\label{lem:potenze-log}
\leavevmode
\begin{enumerate}
  \item Per ogni $c>0$ si ha $n^{c}\to+\infty$, e quindi $n^{-c}\to0$.
  \item Per ogni $x\in\R$ si ha $e^{x}>x$; di conseguenza $\log y<y$ per ogni $y>0$.
  \item $\log n\to+\infty$.
  \item Per ogni $b>0$ si ha $(1+\frac1n)^{b}\to1$.
\end{enumerate}
\end{lemma}

\begin{schema}
\Ipotesi{$c>0$, \ $b>0$, \ $x\in\R$, \ $y>0$.}
\Tesi{(i) $n^{c}\to+\infty$, $n^{-c}\to0$; \ (ii) $e^{x}>x$, $\log y<y$;\\
(iii) $\log n\to+\infty$; \ (iv) $(1+\frac1n)^{b}\to1$.}
\end{schema}

\begin{strategia}
Sono quattro fatti elementari, che servono come mattoni nel teorema successivo. Per (i)
e (iii) basta «invertire» la disuguaglianza $n^{c}>M$ o $\log n>M$ con la funzione
inversa, che è crescente. Per (ii) si confronta $e^{x}$ con la potenza di esponente
intero $e^{\lfloor x\rfloor}$, stimata dal basso con Bernoulli. Per (iv) si intrappola
l'esponente reale $b$ tra $0$ e un intero $k\ge b$, e si usano i carabinieri.
\end{strategia}

\begin{dimostrazione}
(i) Fissato $M>0$, per la monotonia delle potenze \eqref{eq:pot3} si ha $n^{c}>M$ se e
solo se $n>M^{1/c}$. La seconda affermazione segue dalla \eqref{eq:inf5}.

(ii) Se $x\le0$ la tesi è ovvia, perché $e^{x}>0$. Sia $x>0$ e $m=\lfloor x\rfloor\ge0$.
Per la disuguaglianza di Bernoulli, ricordando che $e-1>1$ (\cref{cor:e-tra-2-e-3}),
$e^{m}=\bigl(1+(e-1)\bigr)^{m}\ge1+m(e-1)\ge1+m$ (per $m=0$ entrambi i membri valgono
$1$). Quindi, per la monotonia dell'esponenziale, $e^{x}\ge e^{m}\ge1+m>x$. Per la
seconda affermazione, posto $x=\log y$, si ha $y=e^{\log y}>\log y$.

(iii) Fissato $M>0$, poiché il logaritmo è crescente (\cref{prop:log-monotono}), si ha
$\log n>M$ se e solo se $n>e^{M}$.

(iv) Sia $k\in\N$ con $k\ge b$. Poiché $1+\frac1n>1$, per le \eqref{eq:pot5} si ha
$1\le(1+\frac1n)^{b}\le(1+\frac1n)^{k}$; il terzo membro è il prodotto di $k$ successioni
che tendono a $1$, e tende a $1$. Si conclude con il teorema dei carabinieri.
\end{dimostrazione}

\begin{teorema}[Gerarchia degli infiniti]\label{teo:gerarchia}
Per ogni $b>0$ e ogni $a>1$ valgono i limiti \eqref{eq:gerarchia-limiti}.
\end{teorema}

\begin{schema}
\Ipotesi{$b>0$, \ $a>1$.}
\Tesi{$\dfrac{\log n}{n^{b}}\to0$, \ $\dfrac{n^{b}}{a^{n}}\to0$, \
$\dfrac{a^{n}}{n!}\to0$, \ $\dfrac{n!}{n^{n}}\to0$.}
\end{schema}

\begin{strategia}
Gli ultimi tre limiti sono tutti della forma $a_{n}\to0$ con $a_{n}>0$, e il criterio del
rapporto è lo strumento ideale: il rapporto $a_{n+1}/a_{n}$ di quozienti di potenze e
fattoriali si semplifica in modo spettacolare (ad esempio $\frac{(n+1)!}{n!}=n+1$). Per il
primo limite, invece, il criterio del rapporto non serve (il rapporto tende a $1$), e il
testo di riferimento \cite{marcellini-sbordone} rimanda la dimostrazione a un capitolo
successivo. Possiamo però darla subito, usando la disuguaglianza elementare $\log y<y$
del \cref{lem:potenze-log}: applicata a $y=n^{b/2}$, essa dice che $\log n$ è più piccolo
di un multiplo di $n^{b/2}$, e quindi $\frac{\log n}{n^{b}}$ è più piccolo di un multiplo di
$n^{-b/2}$, che è infinitesimo. Il trucco di usare «metà dell'esponente» è tipico: si
sacrifica un po' della potenza per ottenere una stima semplice.
\end{strategia}

\begin{dimostrazione}
\emph{Primo limite.} Per il \cref{lem:potenze-log}(ii) applicato a $y=n^{b/2}$, e per la
\eqref{eq:log3},
\[
  \frac b2\log n=\log n^{b/2}<n^{b/2},
\]
quindi
\[
  0\le\frac{\log n}{n^{b}}<\frac2b\cdot\frac{n^{b/2}}{n^{b}}=\frac2b\,n^{-b/2}
  \qquad\perogni n\in\N .
\]
Per il \cref{lem:potenze-log}(i) il terzo membro tende a zero, e la tesi segue dal
teorema dei carabinieri.

\emph{Secondo limite.} Poniamo
\begin{equation}\label{eq:gerarchia-2}
  a_{n}=\frac{n^{b}}{a^{n}}; \qquad
  r_{n}=\frac{a_{n+1}}{a_{n}}=\Bigl(\frac{n+1}{n}\Bigr)^{b}\frac1a\to\frac1a<1,
\end{equation}
dove si è usato il \cref{lem:potenze-log}(iv). Dal criterio del rapporto otteniamo che
$n^{b}/a^{n}\to0$.

\emph{Terzo limite.} Poniamo
\begin{equation}\label{eq:gerarchia-3}
  a_{n}=\frac{a^{n}}{n!}; \qquad r_{n}=\frac{a_{n+1}}{a_{n}}=\frac{a^{n+1}}{(n+1)!}\cdot\frac{n!}{a^{n}}=\frac{a}{n+1}\to0 .
\end{equation}
Ancora dal criterio del rapporto segue che $a^{n}/n!\to0$.

\emph{Quarto limite.} Poniamo
\begin{equation}\label{eq:gerarchia-4}
\begin{split}
  a_{n}=\frac{n!}{n^{n}}; \qquad
  r_{n}&=\frac{a_{n+1}}{a_{n}}=\frac{(n+1)!}{(n+1)^{n+1}}\cdot\frac{n^{n}}{n!}\\
  &=\frac{n^{n}}{(n+1)^{n}}=\frac{1}{\bigl(\frac{n+1}{n}\bigr)^{n}}\to\frac1e<1,
\end{split}
\end{equation}
e di nuovo, per il criterio del rapporto, la successione $a_{n}=n!/n^{n}$ converge a zero
per $n\to+\infty$.
\end{dimostrazione}

\begin{osservazione}
Il primo limite vale per i logaritmi in qualsiasi base $c>1$, perché
$\log_{c}n=\log n/\log c$ \eqref{eq:log4}. Il risultato si ricorda con la gerarchia
\[
  \log n\ \ll\ n^{b}\ \ll\ a^{n}\ \ll\ n!\ \ll\ n^{n},
\]
dove $x_{n}\ll y_{n}$ significa $x_{n}/y_{n}\to0$. Si noti che la gerarchia non dipende
dalle costanti: anche $\log n$ è trascurabile rispetto a $n^{0{,}001}$, e $n^{1000}$
rispetto a $1{,}001^{n}$. Ma il «sorpasso» può avvenire molto tardi, come mostra
l'\cref{es:n10-su-1-1n}.
\end{osservazione}

\begin{idea}[Il costo degli algoritmi]
La gerarchia degli infiniti è il linguaggio con cui gli informatici classificano il costo
degli algoritmi. Un problema di dimensione $n$ (ad esempio il numero di celle di una
griglia di calcolo, o di elementi di un modello strutturale) si risolve con un numero di
operazioni che può crescere come $\log n$ (ricerca binaria), come $n$ o $n\log n$
(ordinamento), come $n^{2}$ o $n^{3}$ (eliminazione di Gauss su una matrice piena), come
$2^{n}$ (esplorare tutti i sottoinsiemi) o come $n!$ (provare tutti i percorsi del
problema del commesso viaggiatore). Raddoppiare la potenza del calcolatore permette di
affrontare problemi molto più grandi se il costo è polinomiale, ma aggiunge appena una
unità alla dimensione trattabile se il costo è $2^{n}$. La \cref{fig:infiniti} rende
visibile la differenza.
\end{idea}

\begin{python}[La gerarchia degli infiniti in cifre]
import math

print(f"{'n':>4} {'log n':>8} {'n^2':>8} {'2^n':>12}"
      f" {'n!':>12} {'n^n':>12}")
for n in (5, 10, 20, 50, 100):
    print(f"{n:>4} {math.log(n):8.3f} {n**2:8d} {2.0**n:12.3e} "
          f"{float(math.factorial(n)):12.3e} {float(n)**n:12.3e}")
\end{python}
\begin{risultato}
   n    log n      n^2          2^n           n!          n^n
   5    1.609       25    3.200e+01    1.200e+02    3.125e+03
  10    2.303      100    1.024e+03    3.629e+06    1.000e+10
  20    2.996      400    1.049e+06    2.433e+18    1.049e+26
  50    3.912     2500    1.126e+15    3.041e+64    8.882e+84
 100    4.605    10000    1.268e+30   9.333e+157   1.000e+200
\end{risultato}

\begin{figure}[htb]
\centering
\includegraphics[width=0.85\textwidth]{infiniti}
\caption{Le successioni $\log n$, $n^{2}$, $2^{n}$, $n!$, $n^{n}$ in scala logaritmica
(\texttt{python/fig\_infiniti.py}): le curve si separano sempre di più, ciascuna
diventando trascurabile rispetto alla successiva.}
\label{fig:infiniti}
\end{figure}
